\documentclass[conference]{IEEEtran}

% ---------- Packages ----------
\usepackage{cite}
\usepackage{amsmath,amssymb,amsfonts}
\IEEEoverridecommandlockouts
\usepackage{algorithm}
\usepackage{algorithmic}
\usepackage{graphicx}
\usepackage{booktabs}
\usepackage{siunitx}
\usepackage{textcomp}
\usepackage{xcolor}
\usepackage{hyperref}
\hypersetup{hidelinks}
\usepackage{placeins}
\usepackage{needspace} 
\usepackage{tikz}
\usepackage{quantikz}
\usetikzlibrary{arrows.meta,positioning,fit,calc}
\setlength{\emergencystretch}{1em}

\def\BibTeX{{\rm B\kern-.05em{\sc i\kern-.025em b}\kern-.08em
    T\kern-.1667em\lower.7ex\hbox{E}\kern-.125emX}}

% ---------- Macros (notation) ----------
\newcommand{\Spec}{\mathrm{Spec}}
\newcommand{\Uset}{\mathcal{U}}
\newcommand{\Zunsafe}{Z_{\text{unsafe}}}
\newcommand{\zvec}{\mathbf{z}}
\newcommand{\omegav}{\boldsymbol{\omega}}

\begin{document}

\title{Earliest Failure Localization in Pipelines via Quantum Search}

\author{\IEEEauthorblockN{Alexander Fisher}
\IEEEauthorblockA{\textit{Department of Computer \& Information Science \& Engineering} \\
\textit{University of Florida}\\
Gainesville, FL, USA\\
alexanderfisher@ufl.edu}
\and
\IEEEauthorblockN{Christophe Bobda}
\IEEEauthorblockA{
\textit{Department of Electrical and Computer Engineering}\\
\textit{University of Florida}\\
Gainesville, FL, USA\\
cbobda@ece.ufl.edu
}
}

\maketitle

\IEEEpubid{\makebox[\columnwidth]{979-8-3315-5193-3/26/\$31.00~\textcopyright~2026 IEEE\hfill}%
\hspace{\columnsep}\makebox[\columnwidth]{}} 

\begin{abstract}
Debugging pipelined RTL designs is often delayed by weak localization: a failure may be observed late in a trace even when the first illegal condition occurs earlier and is later masked or overwritten. We formulate earliest-stage localization as a predicate search problem over a fixed environment model. Stage-scoped Boolean (and bounded-temporal) rules induce per-stage violation predicates and an earliest-stage predicate that accepts an assignment only when a stage violates and no earlier stage violates, yielding disjoint attribution by construction. We show how to realize this predicate as a reversible compute--phase--uncompute oracle, and we cast oracle synthesis as a time--space optimization problem on the predicate’s Boolean directed acyclic graph under a hard workspace budget. Our main contribution is a budgeted active-volume compilation pass that emits legal compute/uncompute schedules with peak live intermediates bounded by an ancilla budget, while minimizing a recomputation-aware unit-cost proxy. A grounded case study on an exported 41-node predicate graph demonstrates the induced workspace--recomputation tradeoff and separates exact versus heuristic scheduling behavior in the tight-budget regime.
\end{abstract}

\begin{IEEEkeywords}
hardware verification, staged datapaths, earliest-failure localization, reversible computing, ancilla budgeting, Grover oracle, amplitude amplification
\end{IEEEkeywords}

\IEEEpubidadjcol

\section{Introduction}
Modern hardware pipelines (e.g., CPU front-ends, digital signal processing datapaths, and network accelerators) are engineered as staged, clocked dataflows. When a design violates a safety property, engineers typically need more than a counterexample: they need \emph{localization}---which pipeline stage first becomes inconsistent with the specification. Existing formal methods, including satisfiability modulo theories (SAT/SMT)-based bounded model checking (BMC), IC3/PDR-style model checking, and simulation-guided debug, can produce witnesses, but attributing an \emph{earliest} failing stage can still be expensive when (i) the triggering behaviors are rare, (ii) the symbolic input space is large, or (iii) the search must be repeated across many stage-scoped predicates.\cite{clarke1999modelchecking,clarke2001bmc,bradley2011ic3}

This work frames earliest-failure localization as a \emph{search} problem over a finite Boolean domain. Consider a pipeline partitioned into ordered stages $s \in \{0,\dots,S-1\}$. Let $\zvec \in \{0,1\}^n$ denote a bounded assignment covering whatever bits are treated as symbolic for the run (e.g., primary inputs and/or a bounded unrolling of state/input variables under a chosen horizon). For each stage $s$, a developer supplies (or a tool derives) a set of stage-scoped predicates $\{R_{s,j}(\zvec)\}_{j\in\mathcal{J}_s}$ encoding “this stage violates rule $j$ under assignment $\zvec$.” We summarize stage failure by
\begin{equation}
S_s(\zvec) := \bigvee_{j\in\mathcal{J}_s} R_{s,j}(\zvec),
\end{equation}
and define the earliest-stage predicate
\begin{equation}
G_s(\zvec) := S_s(\zvec)\ \wedge\ \bigwedge_{i<s}\neg S_i(\zvec),
\end{equation}
so that $G_s(\zvec)=1$ means “$\zvec$ is a witness whose first failing stage is exactly $s$.” This definition fixes the meaning of an “earliest-stage hit” independently of whether the backend is classical or quantum.

Importantly, the search variable $\zvec$ is the \emph{global, replayable witness}---the bounded assignment at the pipeline entry and any other symbolic variables admitted by the run model. Stage-local failing values are not searched independently; rather, they are \emph{derived} from $\zvec$ under the compiled stage semantics. Accordingly, a satisfying witness induces both (i) a replayable global assignment and (ii) the concrete stage-local input/signals on which the earliest violating rule is satisfied.

\textbf{From pipelined hardware descriptions to a quantum oracle}
Given a pipelined hardware description (any hardware description language (HDL)), the front-end performs: (i) parse/elaborate across files/modules into an intermediate representation (IR) with resolved symbols and stage boundaries; (ii) bind each rule to concrete stage-local signals/bits; and (iii) lower each rule to a Boolean predicate, compiling purely combinational rules directly and compiling temporal rules by selecting a bound $k$, unrolling to time-indexed bits, then reducing to Boolean form. This produces a Boolean circuit/directed acyclic graph (DAG) computing $G_s(\zvec)$ from $\zvec$, which is then compiled into a reversible \emph{oracle} circuit used by Grover search.\cite{bennett1973logical}

\textbf{Search view (Grover)}
Grover’s algorithm searches the domain $\zvec\in\{0,1\}^n$ for a satisfying assignment to an oracle predicate $f(\zvec)\in\{0,1\}$; here we set $f(\zvec)=G_s(\zvec)$.\cite{grover1996,brassard2002amplification} If there are $K$ satisfying assignments (“hits”) among $N=2^n$ candidates, Grover finds a hit using $O(\sqrt{N/K})$ oracle evaluations.\cite{grover1996,boyer1998tight} In practice, $K$ is typically unknown, so we run a short schedule of iteration counts, measure candidate $\zvec^\ast$ values, and optionally \emph{post-validate classically} by evaluating $G_s(\zvec^\ast)$ and reporting both the triggering rule(s) and the derived stage-local witness.\cite{boyer1998tight} This yields an end-to-end workflow: \emph{(front-end compile)} pipelined design $\rightarrow$ stage predicates $\rightarrow$ earliest predicate $G_s$ $\rightarrow$ \emph{(oracle circuit)} $\rightarrow$ Grover iterations $\rightarrow$ candidate witness $\rightarrow$ classical validation and reporting.

The remainder of the paper focuses on making this oracle construction scalable under realistic qubit constraints by treating reversible oracle synthesis as an explicit time--space optimization problem, rather than assuming that all intermediate values can be stored simultaneously.\cite{bennett1989timespace} We further support this formulation with a concrete compiler-driven case study that measures workspace, recomputation, and depth tradeoffs on exported pipeline stage cones.


\section{Method Overview: From Pipelined Designs to Grover Oracles}
We target \emph{pipelined hardware designs} described by one or more artifacts (HDL, netlist, or IR) together with
developer-specified stage boundaries and stage-scoped correctness rules.
Rules are either combinational predicates over stage-local signals, or bounded temporal rules reduced to Boolean form
via $k$-step unrolling.\cite{clarke2001bmc,BiereEtAl2009BMCHandbook}
We compile the design + rules into Boolean stage predicates over an oracle input vector
$\zvec\in\{0,1\}^n$ (primary inputs plus any nondeterministic/environment bits required by the chosen bound).

\subsection{Stage predicates and earliest-stage localization}
For each stage $i$, compile its rule set into
\[
\begin{aligned}
S_i(\zvec)\;&:=\;\bigvee_{j\in\mathcal{J}_i} R_{i,j}(\zvec),
\qquad
P_s(\zvec)\;:=\;\bigvee_{i<s} S_i(\zvec),\\
G_s(\zvec)\;&:=\;S_s(\zvec)\wedge \neg P_s(\zvec).
\end{aligned}
\]
so $G_s(\zvec)=1$ means: \emph{stage $s$ violates and no earlier stage violates under the same $\zvec$}.
This fixes the meaning of an “earliest-stage witness” independently of backend.
The search variable $\zvec$ is the \emph{global replayable witness}; the concrete stage-local input/signals at the earliest failing stage are derived from $\zvec$ under the compiled stage semantics and can be reported alongside it.

\subsection{Reversible oracle (compute--phase--uncompute)}
Grover requires a phase-marking oracle on the data register:\cite{grover1996,boyer1998tight,brassard2002amplification,nielsenchuang2010}
\[
O_s:\ |\zvec\rangle \mapsto (-1)^{G_s(\zvec)}|\zvec\rangle,
\]
implemented using standard compute--phase--uncompute:\cite{bennett1973logical,Barenco1995ElementaryGates,nielsenchuang2010}
\[
O_s \;=\; U_{G_s}^{\dagger}\,(I\!\otimes\! Z_{\texttt{gflag}})\,U_{G_s},
\]
\[
U_{G_s}:\ |\zvec\rangle|0^{a}\rangle \mapsto
|\zvec\rangle|G_s(\zvec)\rangle|\mathrm{junk}(\zvec)\rangle .
\]
Here $|\mathrm{junk}(\zvec)\rangle$ denotes unavoidable intermediate workspace values from reversible computation;\cite{bennett1973logical,bennett1989timespace}
uncomputing is required so the diffuser acts on a clean $|\zvec\rangle$ space.\cite{brassard2002amplification,nielsenchuang2010}

Concretely, the oracle uses a data register $|\zvec\rangle$, per-rule hit ancillas $a_j$ initialized to $|0\rangle$,
helper bits for OR/prefix aggregation, and a single goal/phase ancilla \texttt{gflag}.
Each rule-hit is computed reversibly by an XOR-write:
\[
|\zvec\rangle|a_j\rangle \mapsto |\zvec\rangle|a_j\oplus R_{s}^{(j)}(\zvec)\rangle,
\]
realized with MCX/Toffoli gates whose controls are \emph{only the literals appearing in the rule} (open controls via $X$-conjugation).\cite{Barenco1995ElementaryGates,nielsenchuang2010}
We reduce rule hits to a stage flag and earlier-stage flags to a prefix flag using De Morgan
(e.g., $a_1\vee\cdots\vee a_m=\neg(\neg a_1\wedge\cdots\wedge \neg a_m)$),
then compute \texttt{gflag}$\leftarrow$\texttt{gflag}$\oplus(\texttt{sflag}\wedge \neg\texttt{prev\_flag})$, apply $Z$, and uncompute all ancillas.

\subsection{Grover loop and readout}
Grover searches $\zvec\in\{0,1\}^n$ with predicate $f(\zvec)=G_s(\zvec)$.\cite{grover1996,boyer1998tight,zalka1999optimal,brassard2002amplification}
If $K$ satisfying assignments exist among $N=2^n$ candidates, a hit is found with
$O(\sqrt{N/K})$ oracle evaluations.\cite{grover1996,boyer1998tight,zalka1999optimal}
Since $K$ is typically unknown, we run a short schedule of iteration counts,
measure a candidate $\zvec^{\ast}$, and optionally post-validate by evaluating $G_s(\zvec^{\ast})$ and reporting the triggered rule(s) together with the derived stage-local witness.\cite{boyer1998tight,brassard2002amplification}

\subsection{Execution regime and relation to prior quantum verification work}
The Grover loop considered here should not be interpreted as a claim of immediate practicality on current noisy intermediate-scale quantum (NISQ) hardware. Rather, we view the present work as a resource-conscious quantum compilation framework for earliest-stage verification search: the main objective is to characterize and optimize oracle workspace/depth tradeoffs under explicit ancilla budgets. These tradeoffs remain relevant under any future execution regime in which larger and more reliable quantum devices make deeper oracle-based search practical, while still providing useful feasibility estimates for small compiled demonstrators on present-day platforms.\cite{preskill2018nisq,endo2018practicalqem,cai2023qem}

This setting also differs from prior quantum formal-verification work, which primarily studies verification of quantum systems, quantum programs, or quantum circuits themselves.\cite{ying2019modelchecking,feng2015qpmc,ying2021quantumcircuits} In contrast, we target classical pipelined hardware designs and use quantum search only as a backend for sampling sparse earliest-stage witnesses.

\section{Contribution: Budgeted Active-Volume Compilation of \texorpdfstring{$O_s$}{Os}}
A naive reversible realization of $U_{G_s}$ may keep many intermediate values live simultaneously, inflating both
workspace qubits and Toffoli/$T$ depth per oracle call. We compile $O_s$ under an explicit ancilla budget by casting
oracle construction as a \emph{budgeted reversible-pebbling} problem on the Boolean DAG for $G_s$.

\paragraph{Model}
Let $\mathcal{C}_{G_s}=(V,E)$ be the Boolean DAG (e.g., AIG) computing $G_s(\zvec)$ from oracle inputs $\zvec$.\cite{abc2010,dagaig2006}
Each node $v\in V$ denotes an intermediate Boolean value (or subexpression) in this computation graph.
A reversible implementation is a dependency-legal schedule
$\pi=(op_1,\ldots,op_T)$, where $op_t\in\{\mathrm{Compute}(v),\mathrm{Uncompute}(v)\}$.\cite{bennett1989timespace}
Let $\mathcal{I}$ denote the set of oracle input variables (the data-register bits corresponding to $\zvec$), and let
$L_t\subseteq V$ be the set of intermediate node-values currently materialized on ancillas after step $t$.
We define the active volume
\[
\mathrm{AV}(\pi):=\max_t |L_t|,
\]
and impose the ancilla budget constraint $\mathrm{AV}(\pi)\le B$.

\paragraph{Legality (reversible pebbling)}
With $\mathrm{fanin}(v)$ the immediate predecessors of $v$, we use the following legality constraints:\cite{MeuliSoekenSchmittMozafariDeMicheli2021ASPDAC,BhattacharjeeSoekenDuttaChattopadhyayDeMicheli2019DATE}
\[
\begin{aligned}
\mathrm{Compute}(v):\quad &\mathrm{fanin}(v)\subseteq (L_t\cup \mathcal{I}),\\
&L_{t+1}=L_t\cup\{v\},\\[2pt]
\mathrm{Uncompute}(v):\quad &v\in L_t,\ \mathrm{fanin}(v)\subseteq (L_t\cup \mathcal{I}),\\
&L_{t+1}=L_t\setminus\{v\}.
\end{aligned}
\]
Thus, a node may be computed or uncomputed only when its fanins are currently available, either as live intermediates or as oracle inputs.

\paragraph{Cost model}
Let $\tau(v)$ denote the cost of a \emph{single} $\mathrm{Compute}(v)$.
Since a clean oracle must also undo any computed node, each time $v$ is computed it is eventually uncomputed as well.
If $r_\pi(v)$ is the number of \emph{recomputations} of $v$ beyond its first computation, then
\[
\mathrm{Cost}(\pi)=2\sum_{v\in V}\tau(v)\big(1+r_\pi(v)\big).
\]
This model makes the time--space tradeoff explicit: tighter ancilla budgets reduce live workspace, but may force additional recomputation.

In the case study, we use a unit-cost instantiation of the model (each \(\mathrm{Compute}\) and \(\mathrm{Uncompute}\) costs 1) to isolate the workspace--recomputation tradeoff; a more detailed Toffoli/$T$ accounting can be substituted without changing the scheduling formulation.

\paragraph{Optimization}
Under this state/cost model, finding a minimum-cost legal schedule can be formulated as a shortest-path problem on a configuration graph.
This is exponential in $|V|$, so we apply exact optimization only to reduced cones or small subgraphs.
For larger $\mathcal{C}_{G_s}=(V,E)$, we run a fixed portfolio of legality-preserving scheduling heuristics, evaluate the emitted schedules under the same cost model, and select
\[
\pi^\star \in \arg\min_{\pi\in\Pi_{\text{port}}}\ \mathrm{Cost}(\pi)
\quad \text{s.t.}\ \mathrm{AV}(\pi)\le B.
\]

\paragraph{Implication}
The compiler enforces $\mathrm{AV}(\pi)\le B$, so oracle workspace is explicitly bounded by the available ancilla budget rather than by the size of the full Boolean DAG.
Any remaining scaling pressure appears as quantified recomputation overhead through $r_\pi(v)$, which directly increases oracle depth.

\section{Case Study and Experimental Setup}

We examine budgeted reversible scheduling on a grounded AIG cone exported from a pipelined RTL benchmark, measuring cost and recomputation as a function of the ancilla budget~$B$.

\subsection{Benchmark and predicate extraction}
\label{subsec:benchmark}

We evaluate the compilation flow on a controlled multi-stage pipelined RTL benchmark with stage-scoped rule annotations. Raw support bits enter at time $t{=}0$; stage~1 computes reusable Boolean facts, stage~2 composes shared intermediates, and stage~3 exposes the final rule-relevant signals. We attach a single Boolean rule at stage~3 and form the earliest-stage predicate
\[
G_s(\zvec)=S_s(\zvec)\ \wedge\ \bigwedge_{i<s}\neg S_i(\zvec),
\]
where $S_i(\zvec)$ is the disjunction of the rules attached to stage $i$.

To isolate reversible-oracle compilation effects, we export a grounded Boolean logic cone for $G_s$ using a SystemVerilog-to-Boolean front-end. The export is a standalone And-Inverter Graph (AIG) cone over a fixed support vector $\zvec$: it provides a root literal for $G_s$, the internal AND nodes with explicit fanins (edge inversions), and the leaf input literals with annotations mapping each literal to a signal/bit/time triple. The internal nodes inline the staged pipeline logic needed to derive the stage-3 rule from $\zvec$, rather than merely restating a shallow final-stage expression.

All results below use one grounded benchmark cone with $|V|=41$ internal AND nodes, constructed to exhibit reconvergence and shared intermediates so that scheduling policies separate under tight budgets.

\subsection{Compilation and scheduling modes}
\label{subsec:comp-modes}

For each exported $G_s$ cone, we compile a reversible schedule under an active-volume budget $B$ (inputs/constants are always available and excluded from the live set). We compare four schedule-generation modes. \emph{Store-all} computes all internal nodes once and keeps them live until the final cleanup, giving a full-workspace baseline with zero recomputation. \emph{Oldest-live} enforces the budget by evicting (uncomputing) the longest-lived intermediate when space is needed, serving as a weak heuristic baseline. \emph{Portfolio-best} uses a shared recursive ``make-the-root-live'' scheduling engine and runs three greedy eviction policies (smallest-fanout, depth-aware, recompute-aware), selecting the lowest-cost feasible schedule produced. \emph{Exact-optimal} treats scheduling as a shortest-path problem over legal live-set configurations (live mask plus a root-used flag) and uses \emph{Dijkstra's algorithm} to certify minimum-cost schedules where tractable.\cite{dijkstra1959note}

\subsection{Measured quantities}
For each exported cone, budget $B$, and schedule mode, we record:
\begin{itemize}
    \item peak live intermediates (active volume) $\mathrm{AV}(\pi)$;
    \item total compute operations and total uncompute operations;
    \item total recomputations $\sum_{v\in V}(\mathrm{compute\_count}(v)-1)$;
    \item total schedule cost under the unit-cost model,
    $\mathrm{Cost}(\pi)=\#\mathrm{Compute}+\#\mathrm{Uncompute}$;
\end{itemize}

These measurements directly quantify the workspace--recomputation tradeoff induced by ancilla-bounded scheduling on a fixed predicate DAG: as $B$ decreases, fewer intermediates can remain live simultaneously, which may force recomputation and increase total cost.


\section{Results and Discussion}
\begin{figure}[t]
  \centering
  \includegraphics[width=\columnwidth]{figures/cost_vs_budget.pdf}
  \caption{Budgeted reversible compilation cost versus ancilla budget $B$ on the 41-node grounded cone. Black diamonds denote \emph{certified} exact-optimal points (connected by a solid segment). The dotted black curve and gray ribbon indicate an \emph{estimated} exact trend and range where exhaustive certification was not performed. Orange: oldest-live heuristic. Blue: portfolio-best heuristic. The vertical line marks the full-workspace baseline at $B=|V|=41$.}
  \label{fig:cost-vs-budget}
\end{figure}

\subsection{Budget--cost tradeoff}
\label{subsec:cost-tradeoff}

Fig.~\ref{fig:cost-vs-budget} shows the expected workspace--cost tradeoff induced by ancilla-bounded scheduling on a fixed predicate DAG. At tight budgets, fewer intermediates can remain live simultaneously, forcing additional recomputation and increasing total cost. The certified exact-optimal points illustrate that feasible schedules exist at substantially smaller budgets than the heuristics discover: at $B\in\{7,8,9\}$, exact schedules are feasible with costs $222$, $184$, and $156$, respectively, while the heuristics remain infeasible until larger budgets.

Once the heuristic methods become feasible, portfolio-best is never worse than the weak oldest-live baseline on the feasible overlap, and provides a modest improvement near the tight end of that region. As $B$ increases, all methods converge toward the full-workspace store-all endpoint at $B=|V|=41$, where recomputation vanishes and total cost reaches $82$. For intermediate budgets where exact optimization is not exhaustively certified, Fig.~\ref{fig:cost-vs-budget} distinguishes certified exact points from an estimated exact trend/range to reflect the computational limits of configuration-graph search on the 41-node cone.

\FloatBarrier
\begin{figure}[t]
  \centering
  \includegraphics[width=\columnwidth]{figures/recomputations_vs_budget.pdf}
  \caption{Total recomputations versus ancilla budget $B$ on the 41-node grounded cone. Black diamonds denote \emph{certified} exact-optimal points (connected by a solid segment). The dotted black curve and gray ribbon indicate an \emph{estimated} exact trend and range beyond the certified segment. Orange: oldest-live. Blue: portfolio-best. The vertical line marks $B=|V|=41$, where recomputation is zero.}
  \label{fig:recomp-vs-budget}
\end{figure}

\subsection{Budget--recomputation tradeoff}
\label{subsec:recomp-tradeoff}

Fig.~\ref{fig:recomp-vs-budget} reports the total number of recomputations incurred by each policy as a function of the ancilla budget. In the low-budget regime, eviction is unavoidable and schedules repeatedly re-materialize intermediate values to re-establish required fanins. The certified exact-optimal points quantify how quickly this overhead can shrink with modest additional workspace: recomputations fall from $70$ at $B{=}7$ to $51$ at $B{=}8$ and $37$ at $B{=}9$, and approach zero as $B$ increases toward the full-workspace endpoint $B{=}|V|{=}41$.

On the feasible overlap, portfolio-best consistently matches or improves upon the weak oldest-live baseline, with the largest differences appearing near the tight end where evictions are frequent; at larger budgets the heuristics converge as recomputation becomes rare. Finally, the close visual agreement between Fig.~\ref{fig:cost-vs-budget} and Fig.~\ref{fig:recomp-vs-budget} is expected under our unit-cost model: total cost is an affine function of recomputation count (each additional recomputed node contributes one extra compute and one extra uncompute), so reducing recomputation directly reduces total cost.

\section{Conclusion}
We formulated earliest-stage failure localization as a disjoint predicate family $G_s(\zvec)$, so each violating assignment is attributed to its first failing stage by construction and can be marked by a Grover-compatible compute--phase--uncompute oracle. We then cast reversible oracle synthesis for $G_s$ as a time--space optimization problem on the predicate graph, with an active-volume budget $B$ over live intermediates and a recomputation-aware cost model. On a grounded 41-node cone extracted from a staged RTL benchmark, the measured cost and recomputation curves exhibit the expected workspace--recomputation tradeoff, and certified exact-optimal points show feasible low-budget schedules that the tested heuristics reach only at larger budgets. Overall, earliest-stage localization can be compiled into stage-guarded oracles whose feasibility and cost can be treated as explicit optimization targets under workspace constraints.

\bibliographystyle{IEEEtran}
\bibliography{refs}

\end{document}
